\begin{abstract}
Large language models are increasingly deployed with confidential information embedded in system prompts, yet the effectiveness of concealment instructions remains poorly understood.
We test whether LLMs leak information they are instructed to conceal at higher rates than information presented neutrally---a ``\streisand'' for language models.
Across 1,680 controlled trials on \gptfour and \gptmini, we vary a single factor---whether the system prompt frames information as confidential or neutral---while holding the information content, probe queries, and all other variables constant.
Contrary to our hypothesis, concealment instructions dramatically \emph{reduce} leakage: from 58.8\% to 2.3\% attack success rate under strict detection ($\chi^2 = 630.5$, $p < 10^{-139}$, Cram\'er's $V = 0.61$).
Paired analysis reveals zero cases where concealment-framed information leaked but its neutral counterpart did not.
The 19 concealment breaches we observe trace to a single attack vector---system prompt translation---which bypasses concealment by reframing disclosure as a language task.
Our results suggest that explicit concealment instructions are highly effective for current frontier models, though they remain vulnerable to specific prompt injection techniques that operate outside the disclosure-decision pathway.
\end{abstract}
